Este artigo quantifica o ganho de produtividade total dos fatores (PTF) necessário para neutralizar, no curto prazo, a perda de produto associada a tetos alternativos para a jornada semanal no Brasil. Combinamos momentos da PNAD Contínua de 2024T4 com um modelo estático de equilíbrio parcial no qual trabalho formal e informal são substitutos imperfeitos, a eficiência horária depende da jornada e as cunhas de formalização governam a composição do emprego. Para representar as reformas de 44 para 40 e 36 horas, limitamos previamente as horas habituais formais a 44h, como hipótese de mensuração da base. A alternativa de 40 horas requer $1{,}53$--$1{,}57\%$ de PTF. A passagem para 36 horas requer entre $6{,}52\%$ e $8{,}38\%$, conforme a hipótese sobre eficiência abaixo de 40 horas; no mesmo contrafactual, a informalidade aumenta entre $2{,}28$ e $2{,}92$ pontos percentuais. Decomposições mostram que a contração física de horas formais explica a maior parcela negativa da variação do produto. A comparação com a trajetória brasileira de PTF e os exercícios de bem-estar sugerem uma diferença quantitativa importante entre reformas moderadas e o teto de 36 horas. As magnitudes dependem da fadiga e das horas iniciais: preservar a cauda habitual acima de 44h reduz o requisito em 40 horas para cerca de $0{,}07\%$. Os resultados motivam a avaliação de trajetórias graduais e de instrumentos de transição, sem constituir uma estimativa causal dos efeitos de uma proposta legislativa específica.
